\section{Operadores de reparación}
\label{sec:reparacion}

Los cinco operadores tienen la firma $\textit{rep}(x,\textit{ruido})$ y reinsertan los
clientes de $\Uns(x)$ uno a uno. Todos son heurísticas constructivas de inserción: en cada
paso eligen un cliente y una posición según su criterio, ejecutan la inserción y repiten. Son
instancias de una misma plantilla (Sección~\ref{sec:cache}) y comparten las siguientes reglas.

\begin{itemize}
\item \textbf{Rutas admisibles.} Una ruta admite al cliente $u$ si
$\textit{carga}+\mu_u\le Q$. Las rutas vacías que haya dejado la destrucción son rutas
admisibles como cualquier otra.
\item \textbf{Costo de inserción.} Es el incremento de distancia $\Delta$ del
Algoritmo~\ref{alg:evalins}, exacto o perturbado con ruido (Sección~\ref{sec:ruido}). Insertar
en una ruta vacía cuesta $2\,d_{0u}$ y \emph{no} incluye la penalización de $10^4$ por
vehículo: la reparación razona solo en distancia. El control del número de vehículos recae en
la estructura en dos fases de la búsqueda y en la función de costo, a través del criterio de
aceptación.
\item \textbf{Flota fija.} La reparación no abre rutas nuevas. Si ningún cliente pendiente
tiene una inserción factible en las rutas existentes, termina y deja esos clientes en $\Uns$.
La candidata resultante es incompleta y el bucle común la evalúa con la penalización de $10^6$
por cliente sin asignar (Sección~\ref{sec:fases}).
\end{itemize}

\subsection{Formulación directa y su costo}

Sea $U=|\Uns|$ el número de clientes por insertar: $U=q$ tras una destrucción en la fase~2, y
$q$ más los clientes del banco en la fase~1. La formulación directa evalúa en cada paso todas
las posiciones de todas las rutas para todos los clientes pendientes: $\bigO(U\,n)$
evaluaciones por paso y $\bigO(U^2 n)$ en total. Aun con $q\le 60$, esto supone del orden de
$10^6$ evaluaciones de inserción por iteración para $n=800$, demasiado para un presupuesto de
decenas de miles de iteraciones por corrida.

\subsection{Inserción con caché por ruta}
\label{sec:cache}

La observación clave es que una inserción modifica \emph{una sola ruta}. Las evaluaciones de
todos los demás pares (cliente, ruta) siguen siendo válidas. El
Algoritmo~\ref{alg:cache} (\cod{cachedInsertion}) mantiene una matriz
$M[u][r]$ con un resumen de la ruta $r$ para el cliente pendiente $u$:
\[
M[u][r]=(\textit{mejor},\ \textit{segundo},\ \textit{pos}),
\]
el menor y el segundo menor costo de inserción factible de $u$ en $r$ y la posición del
primero ($\infty$ si no existen). Calcular una entrada cuesta $\bigO(L_r)$ gracias a la
evaluación en tiempo constante. Tras cada inserción en la ruta $r^\star$ solo se recalcula la
columna $r^\star$.

Sobre cada fila se define un \emph{agregador} $G[u]$ que la resume en lo que el criterio
necesita: la ruta de destino y una \emph{clave} de prioridad. Los cinco operadores difieren
únicamente en el agregador (Sección~\ref{sec:agregadores}).

\begin{algorithm}[H]
\caption{\fn{InserciónConCaché}$\langle G\rangle$ (\cod{cachedInsertion<Agg>})}
\label{alg:cache}
\Entrada{solución $x$ con $\Uns\neq\emptyset$; bandera \textit{ruido}; tipo de agregador $G$}
$\nu\gets 0.025\,d_{\max}$ si \textit{ruido}; si no, $\nu\gets 0$ \Com*[r]{amplitud del ruido}
\ParaCada{$u\in\Uns$}{
  \lParaCada{ruta $r$}{$M[u][r]\gets$ \fn{Entrada}($u$, $r$, $\nu$)}
  $G[u].\fn{reconstruir}(M[u])$\;
}
\Mientras{$\Uns\neq\emptyset$}{
  $u^\star\gets\arg\max\{G[u].\textit{clave}:\ u\in\Uns,\ G[u]\text{ válido}\}$ \Com*[r]{primer máximo}
  \lSi(\Com*[f]{flota fija: los restantes quedan sin asignar}){no existe $u^\star$}{\Salir}
  $r^\star\gets G[u^\star].\textit{ruta}$\;
  insertar $u^\star$ en $r^\star$ en la posición $M[u^\star][r^\star].\textit{pos}$;\quad \fn{RecalcularRuta}($r^\star$)\;
  quitar $u^\star$ de $\Uns$, de $M$ y de $G$ \Com*[r]{en $\bigO(1)$, intercambiando con el último}
  \ParaCada(\Com*[f]{solo cambió la columna $r^\star$}){$u\in\Uns$}{
    $e_{\mathrm{ant}}\gets M[u][r^\star]$;\quad $M[u][r^\star]\gets$ \fn{Entrada}($u$, $r^\star$, $\nu$)\;
    \lSi{$M[u][r^\star]\neq e_{\mathrm{ant}}$}{$G[u].\fn{actualizar}(M[u],\ r^\star,\ e_{\mathrm{ant}})$}
  }
}
actualizar $K(x)$ y $D(x)$\;
\end{algorithm}

\begin{algorithm}[H]
\caption{\fn{Entrada} (\cod{computeEntry})}
\label{alg:entrada}
\Entrada{cliente $u$, ruta $r$ con secuencia $\pi$, amplitud de ruido $\nu$}
\Salida{$e=(\textit{mejor},\ \textit{segundo},\ \textit{pos})$}
$e\gets(\infty,\infty,-)$\;
\lSi{$\textit{carga}_r+\mu_u>Q$}{\Devolver $e$}
\Para{$k\gets 0$ \Hasta $L-2$}{
  $(\textit{factible},\Delta)\gets$ \fn{EvaluarInserción}($u$, $\pi$, $k$)\;
  \lSi{$\neg\,\textit{factible}$}{\Continuar}
  \lSi{$\nu>0$}{$\Delta\gets\max\bigl(0,\ \Delta+\xi\bigr)$ con $\xi\sim\unif{-\nu}{\nu}$}
  \uSi{$\Delta<e.\textit{mejor}$}{
    $e.\textit{segundo}\gets e.\textit{mejor}$;\quad $e.\textit{mejor}\gets\Delta$;\quad $e.\textit{pos}\gets k+1$\;
  }
  \SiNoSi{$\Delta<e.\textit{segundo}$}{
    $e.\textit{segundo}\gets\Delta$\;
  }
}
\Devolver $e$\;
\end{algorithm}

\paragraph{Costo.} La inicialización cuesta $\bigO(U\,n)$ evaluaciones. Cada paso posterior
recalcula una columna, $\bigO(U\,\bar{L})$, y elige el máximo en $\bigO(U)$. El total es
\[
\bigO\bigl(U\,n+U^{2}\,\bar{L}\bigr)\quad\text{frente a}\quad\bigO\bigl(U^{2}\,n\bigr),
\]
una reducción por un factor del orden del número de rutas $m=n/\bar{L}$ en el término
dominante. A esto se suma la actualización de los agregadores, que se analiza a continuación.

\subsection{Ruido en los costos de inserción}
\label{sec:ruido}

Cuando la bandera \emph{ruido} está activa (lo está en la mitad de las iteraciones,
Sección~\ref{sec:catalogo}), cada costo de inserción factible se perturba con un término
aditivo, como en~\cite{ropke2006}:
\begin{equation}
\tilde{\Delta}=\max\bigl(0,\ \Delta+\xi\bigr),\qquad \xi\sim\unif{-\nu}{\nu},\qquad
\nu=0.025\,d_{\max},
\label{eq:ruido}
\end{equation}
donde $d_{\max}$ es la mayor distancia entre dos nodos de la instancia y $0.025$ es la
constante \cod{NOISE\_RATE}. La amplitud es absoluta, proporcional a la escala geográfica de
la instancia y no al costo de cada inserción: altera el orden entre inserciones cuya
diferencia es menor que $2\nu$ y respeta las decisiones claras, de modo que dos reparaciones
de la misma solución parcial pueden producir resultados distintos. Es el mecanismo de
diversificación de la reparación, y se aplica por igual a los cinco operadores.

El ruido se sortea al calcular cada entrada (Algoritmo~\ref{alg:entrada}), con un valor
independiente por posición, y actúa \emph{antes} de seleccionar el mejor y el segundo mejor
costo de la ruta. En consecuencia afecta a las tres decisiones de la reparación: en qué
posición de una ruta entraría un cliente, a qué ruta va y en qué orden se insertan los
clientes. La entrada perturbada se conserva en la caché mientras su ruta no cambie y se
vuelve a sortear cuando la columna se recalcula. La factibilidad no se perturba: una
inserción elegida con ruido es siempre factible, y la distancia de la solución se calcula con
los costos reales.

\subsection{Agregadores y actualización perezosa}
\label{sec:agregadores}

Reconstruir un agregador exige recorrer su fila, $\bigO(m)$. Hacerlo para todos los clientes
en cada paso añadiría $\bigO(U^2 m)$. Cada agregador implementa por ello una operación
\fn{actualizar} que, conocida la entrada anterior y la nueva de la única columna modificada,
decide en $\bigO(1)$ si el resumen puede haber cambiado y solo en ese caso reconstruye. Además,
\fn{actualizar} ni siquiera se invoca cuando la entrada recalculada es idéntica a la anterior.
Las condiciones de salto están escritas de modo que el resultado coincide siempre con el de
una reconstrucción completa, incluidos los empates.

\paragraph{Inserción voraz (\cod{GreedyAgg}).} Resume la fila en el menor costo y su ruta:
\[
c_u=\min_r M[u][r].\textit{mejor},\qquad \textit{clave}=-c_u .
\]
Se inserta primero el cliente cuya mejor inserción es la más barata. Actualización: si la
columna modificada es la ruta óptima actual, se reconstruye; si es otra, basta comparar su
nuevo costo con $c_u$ (y, en caso de empate, preferir la ruta de menor índice, que es lo que
elegiría la reconstrucción). Es el criterio más miope: tiende a agotar pronto las inserciones
baratas y a dejar para el final a los clientes difíciles.

\paragraph{Arrepentimiento 2 (\cod{Regret2Agg}).} Considera los dos menores costos de
inserción del cliente sobre \emph{todas las posiciones de todas las rutas}, $c^{(1)}_u\le
c^{(2)}_u$ (pueden pertenecer a la misma ruta; por eso cada entrada de $M$ guarda también el
segundo mejor costo de la ruta):
\[
\textit{clave}=c^{(2)}_u-c^{(1)}_u,\qquad \textit{clave}=10^{9}\ \text{si solo existe una inserción factible.}
\]
Se inserta primero el cliente que más perdería si no obtuviera su mejor posición; un cliente
con una única posición factible tiene prioridad absoluta. Actualización: si la columna
modificada no es la ruta óptima y tanto su mejor costo anterior como el nuevo superan
$c^{(2)}_u$, ni $c^{(1)}_u$ ni $c^{(2)}_u$ pueden cambiar y se omite la reconstrucción.

\paragraph{Arrepentimiento $k$ sobre rutas (\cod{RegretAgg<K>}).} Los tres operadores
restantes comparten un agregador genérico, parametrizado por $k$. Los valores usados son
$k=3$ (\cod{regret3Insertion}), $k=4$ (\cod{regret4Insertion}) y $k=m$, el número de rutas de
la solución (\cod{regretMInsertion}, que en el código corresponde a \cod{K = 0}). Es el
arrepentimiento de~\cite{ropke2006}: considera el mejor costo \emph{por ruta} y ordena las
rutas factibles, $c^{[1]}_u\le c^{[2]}_u\le c^{[3]}_u\le\dots$ Con
$h=\min(k,\text{número de rutas factibles})$,
\begin{equation}
\textit{clave}=\sum_{l=2}^{h}\bigl(c^{[l]}_u-c^{[1]}_u\bigr)+10^{4}\,(k-h).
\label{eq:regretk}
\end{equation}
El primer término mide cuánto se perdería en total si el cliente tuviera que conformarse con
su segunda, tercera, \dots, $k$-ésima mejor ruta. El segundo prioriza a los clientes con menos
de $k$ rutas factibles, con una penalización de $10^4$ por cada ruta que falta. La ruta de
destino es la de menor costo (la de menor índice en caso de empate).

El agregador conserva, además de la clave y la ruta, el $k$-ésimo menor costo $c^{[k]}_u$
(o $\infty$ si hay menos de $k$ rutas factibles). Actualización: si el mejor costo de la
columna modificada no cambió, o si tanto el anterior como el nuevo superan $c^{[k]}_u$, la
ruta estaba y sigue estando fuera de las $k$ mejores y la clave no cambia; en otro caso se
reconstruye, con un ordenamiento $\bigO(m\log m)$ sobre un búfer reutilizado para evitar
reservas de memoria. Para $k=3$ y $k=4$ el salto se aplica siempre que la columna modificada
no esté entre las $k$ mejores rutas del cliente. Para $k=m$ ninguna ruta factible queda fuera,
de modo que el agregador se reconstruye siempre que cambia el mejor costo de la columna: el
arrepentimiento sobre todas las rutas es el operador de reparación más costoso.

\paragraph{Diferencias entre los arrepentimientos.} El arrepentimiento 2 se mide sobre
posiciones y los demás sobre rutas. El primero es sensible a la diferencia entre las dos
mejores posiciones aunque estén en la misma ruta, y por tanto prioriza a los clientes cuya
colocación \emph{exacta} importa; los segundos miden cuántas rutas alternativas tiene
realmente el cliente y priorizan a los que dependen de pocas rutas. Entre estos, $k$ regula el
horizonte de anticipación: con $k=3$ o $k=4$ solo cuentan las mejores alternativas, mientras
que con $k=m$ el término $10^{4}(m-h)$ suele dominar la clave y el operador inserta primero,
en esencia, a los clientes que caben en \emph{menos} rutas, usando la suma de diferencias
para desempatar. Frente a la inserción voraz, todos anticipan qué clientes quedarán sin
alternativas, lo que es especialmente valioso con ventanas de tiempo estrechas y, sobre todo,
con flota fija: en la fase~1 un cliente que se queda sin ruta factible permanece en el banco
(Sección~\ref{sec:fases}).

\subsection{Resumen de las optimizaciones}

La Tabla~\ref{tab:optim} reúne las optimizaciones descritas en esta sección y en la anterior.
Todas son \emph{exactas}: cambian el costo computacional, no la decisión que toma el operador.
En las reparaciones con ruido, la perturbación forma parte de la definición del operador (se
sortea una vez por entrada calculada) y la caché no introduce ninguna aproximación adicional.

\begin{table}[!htb]
\centering
\small
\caption{Optimizaciones algorítmicas de los operadores. $U$: clientes por insertar;
$\bar{L}$: longitud media de ruta; $m$: número de rutas.}
\label{tab:optim}
\begin{tabularx}{\textwidth}{@{}l L l l@{}}
\toprule
Componente & Técnica & Costo directo & Costo optimizado \\
\midrule
Factibilidad de inserción & Llegadas, esperas y holguras por ruta & $\bigO(\bar{L})$ & $\bigO(1)$ \\
Arcos imposibles & Matriz de alcanzabilidad precalculada & -- & filtro $\bigO(1)$ \\
Extracción de peores & Lista ordenada incremental; solo se actualizan dos vecinos & $\bigO(q\,n\log n)$ & $\bigO(n\log n+q\,n)^{\dagger}$ \\
Extracción de Shaw & Orden de afinidad precalculado por instancia & $\bigO(q\,n\log n)$ & $\bigO(q\,n)^{\ddagger}$ \\
Voraz y arrepentimientos & Caché por (cliente, ruta); solo se recalcula la ruta modificada & $\bigO(U^{2}n)$ & $\bigO(Un+U^{2}\bar{L})$ \\
Agregadores & Actualización perezosa con prueba $\bigO(1)$ & $\bigO(m)$ por cliente & $\bigO(1)$ habitual$^{\S}$ \\
Borrado de pendientes & Intercambio con el último elemento & $\bigO(U)$ & $\bigO(1)$ \\
\bottomrule
\end{tabularx}

\vspace{2pt}
\raggedright\footnotesize
$^{\dagger}$ El término $q\,n$ corresponde a desplazamientos de memoria contigua, con constante muy pequeña.\\
$^{\ddagger}$ Cota de peor caso del recorrido; en la práctica es mucho menor por el sesgo hacia rangos bajos. Se añade un precálculo único de $\bigO(n^{2}\log n)$.\\
$^{\S}$ Salvo en el arrepentimiento sobre todas las rutas, que reconstruye su agregador en cada cambio.
\end{table}
